\section{Experimental design}
\label{sec:design}

\subsection{Conditions}

\begin{table}[t]
\centering
\small
\caption{Planned experimental conditions. Status reflects the repository state at
the time of writing; no condition has produced a research result.}
\label{tab:conditions}
\begin{tabular}{@{}lll@{}}
\toprule
ID & Condition & Status \\
\midrule
E0 & Qwen3.5-0.8B instruct reference; base reference where available & Planned \\
E1 & \cyberlite{} LoRA (bf16) & Planned \\
E2 & \cyberlite{} QLoRA (4-bit) & Planned \\
E3 & \cyberlite{} offline-distilled student & Planned \\
E4 & E1/E2 with BM25 retrieval (RAG) & Planned \\
E5 & Quantized inference: bnb 4/8-bit, TorchAO, GGUF & Planned \\
E6 & CPU-only deployment profile & Planned \\
E7 & Larger general-purpose control models (optional) & Optional \\
\bottomrule
\end{tabular}
\end{table}

E0 establishes the unadapted baseline under identical prompts. E1 and E2 test
the efficiency hypothesis at matched data and steps. E3 tests whether a larger
teacher adds value. E4 tests retrieval as a substitute for parametric knowledge.
E5 tests deployment quantization, and E6 tests the CPU constraint. E7 is not
required for the core claims and is only included when compute permits; mixing
different base models requires reporting tokenizer, prompt, and license
differences.

\subsection{Hardware profiles}

Experiments are grouped by target VRAM class: 4, 6, 8, 12, and 16 GB, plus CPU.
Each run records the exact GPU/CPU model, driver, CUDA, PyTorch, Transformers,
attention backend, dtype, and quantization. A configuration is reported as
feasible only if it completes with measured peak memory below the target and
without out-of-memory errors. Planning estimates are not feasibility evidence.

\subsection{Data and evaluation separation}

Training data comes only from the prepared corpus. Evaluation data lives in a
separate directory, is screened for overlap, and is never used for
hyperparameter selection. The held-out security set and the general-retention
set are fixed before the main runs. Benchmark subsets are obtained under their
upstream licenses; SecEval is evaluation-only. Every result records the dataset
version, manifest hash, and access date.

\subsection{Prompting and decoding}

All conditions use the same prompt templates for a given task. Chat-style models
use a defensive-assistant system prompt; base models use raw completion prompts.
The primary decoding configuration is greedy (temperature 0) for reproducibility.
A secondary sample-based run may be reported for autocomplete, with fixed seeds
and reported variance. Maximum new tokens are fixed per task and reported.

\subsection{Statistical analysis}

For accuracy, we report point estimates and 95\% percentile bootstrap confidence
intervals \citep{efron1979bootstrap}. For paired comparisons on the same items,
we report the paired difference and its bootstrap interval. Classification
results include macro F1 and false-positive rate. Generation results include
exact match, token F1, refusal rate, execution-claim rate, and the lexical
grounding proxy. Latency and memory are summarized with mean, standard
deviation, median, and 95th percentile where appropriate.

We do not use a single composite score to declare a winner. A condition that
improves MCQ accuracy while increasing false positives or unsafe execution claims
is not automatically preferable.

\subsection{Stopping and exclusion rules}

\begin{itemize}[leftmargin=1.4em]
  \item A run that fails to produce a finite loss or save a checkpoint is
  reported as failed, not silently excluded.
  \item Out-of-memory runs are reported with the attempted configuration; the
  next lower sequence length or rank is documented.
  \item Evaluation items with malformed gold labels are corrected before the run
  and the correction is logged.
  \item No evaluation item is removed because a model answers it incorrectly.
  \item Hyperparameters are selected on validation data, not test data.
\end{itemize}

\subsection{Compute and cost accounting}

For each run we record wall-clock duration, peak VRAM, process RSS, and
tokens/second. Cost is computed only from user-supplied rates. Two cost views are
reported when possible: a compute rate (for rented hardware) and an energy
estimate (watts $\times$ hours $\times$ price per kWh). We avoid presenting a
single dollar figure without its assumptions.
